\section{Related Work}
\label{sec:related}

\paragraph{Connectome-constrained computation.}
Large-scale fly connectomes support models constrained by measured neuronal
wiring~\cite{dorkenwald2024,malecns2026}. Lappalainen et
al.~\cite{lappalainen2024} combine connectomic constraints with task training
to predict neural responses across the visual system. Structure-only models
investigate visual computation under fixed wiring and wiring-cost
comparisons~\cite{structure2026}. Beiran and Litwin-Kumar~\cite{beiran2025} characterize dynamical
degeneracy under fixed connectomic weights and show how partial activity
recordings can resolve it. Our question concerns the source information
left by an intervention in a fixed computational model.

\paragraph{Null models and structural attribution.}
Degree-preserving swaps and motif counts are established tools for network
analysis~\cite{maslov2002,milo2002}; weighted-network randomization further
shows that degree and strength are different constraints~\cite{milisav2025}.
Dhiman~\cite{dhiman2026} studies staged controls for initialization, degree,
and early optimization in connectome-constrained networks. The closest
mechanistic precedent is Park~\cite{park2026}, who shows that standard
randomizations introduce sensory--motor shortcuts in embodied agents,
and evaluates boundary-preserving controls, graded shortcut doses, and sham
interventions. These controls address task-relevant routes. Our audit resolves source
access inside the executed propagator, including separate-relay routes
that survive a zero common-relay count. It uses standard linear transport
algebra to diagnose the functional meaning of the structural comparison.
Zhang et al.~\cite{zhang2025motifs} analyze how coherent and incoherent
motifs affect representational capacity and stability; our question is
whether removing a declared motif removes the inputs required by a task.

\paragraph{Graph communication and rewiring.}
Curvature-based analysis~\cite{curvature2022} and joint width--depth--topology
analysis~\cite{oversquashing2023} relate graph structure to information
transport. Spectral pruning demonstrates that removing edges can improve
communication and generalization~\cite{pruning2024}, while virtual-node
rewiring creates efficient long-range routes~\cite{virtual2024}.
GraphGPS combines local and global communication~\cite{graphgps2022};
oversmoothing itself depends on the model's parameter
regime~\cite{phase2024}. We fix the two-step propagator and evaluate what task information its
structural interventions preserve, complementing methods that alter
communication to improve learning.

\paragraph{Evaluating explanations through interventions.}
GNNExplainer uses subgraphs and features to explain learned
predictions~\cite{gnnexplainer2019}. ROAR shows that deletion-based image
attribution can be confounded by distribution shift and evaluates removal
with retraining~\cite{roar2019}. RegExplainer addresses explanation/input
shift and the different decision structure of graph
regression~\cite{regexplainer2024}. Our readout is trained separately for
each graph condition, so the outcome concerns available task information
under that condition, rather than the acute sensitivity of one frozen
trained model. The additional issue here is graph compensation: preserving
degree while removing an edge can introduce a replacement computation.
